\section{总结与延伸}

这堂 Overview 从 attention 的信息选择机制一路走到 agent 的安全边界。表面上主题很多，底层可以压成三次边界扩张：第一，Transformer 让每个 token 直接读取相关上下文；第二，LLM 用规模、数据和 post-training 扩大可处理任务；第三，agent 在模型外增加 memory、tool、state、communication 与 environment feedback。每次扩张都带来新能力，也把验证从单个输出推向完整系统。

\subsection{十个核心结论}

\begin{importantbox}{从模型机制到系统工程}
\begin{enumerate}
  \item Attention 用 query--key matching 选择 value，multi-head 提供多个表示子空间；
  \item Transformer 以 $O(n^2)$ 全局配对换取短信息路径和训练并行性；
  \item LLM 仍以 next-token objective 为基础，instruction following 与 safety 来自额外训练和系统约束；
  \item Emergence curve 既可能反映 scale，也可能被 metric threshold 放大；
  \item RLHF、DPO、MoE 与 data quality 是不同层次的设计，不能压成“模型更大”；
  \item 外部 memory、RAG、continual learning 与 model editing 改变的对象不同；
  \item CoT、ToT 与 Socratic decomposition 增加中间计算，但必须配合 verifier；
  \item Agent 是 observation--action--feedback loop，不是一次模型调用；
  \item Multi-agent 的主要难点是 state、protocol、verification 与 recovery，而不是并行启动；
  \item LLM OS 是组件地图，真正部署还需要 permission、isolation、budget、fallback 与 sandbox。
\end{enumerate}
\end{importantbox}

\subsection{一张表串起整堂课}

\begin{center}
\begin{longtable}{L{0.17\textwidth} L{0.22\textwidth} L{0.24\textwidth} L{0.25\textwidth}}
\toprule
层次 & 新增能力 & 新增失败模式 & 必要验证 \\
\midrule
\endhead
Embedding & 共享连续表示 & 相似度不等于事实关系 & 近邻、偏差与分布测试 \\
Attention & 按内容读取上下文 & 权重难以直接解释、长序列二次成本 & ablation、causal test、效率测试 \\
LLM scale & 更广任务覆盖与 few-shot behavior & 成本、污染、不可预测能力 & 多维 benchmark、数据审计 \\
Preference tuning & 更符合人类或组织偏好 & reward hacking、偏好偏差 & held-out preference、安全评估 \\
External memory & 可更新事实与个性化 & 错误检索、越权、过期 memory & 检索、证据与权限检查 \\
Reasoning scaffold & 更多 test-time compute 与可见步骤 & 自洽但错误、搜索成本 & 工具、验证与反事实测试 \\
Agent loop & 执行长期外部任务 & side effect、loop、plan divergence & 状态日志、预算、回退与恢复 \\
Multi-agent & 并行和 specialization & 协议歧义、冲突、重复副作用 & 结构化消息、版本和独立复核 \\
\bottomrule
\end{longtable}
\end{center}

\subsection{课堂结尾真正留下的问题}

讲者的 closing discussion 把 agent 未来集中在 error correction、security、permission 与 sandbox。这里还有几条值得继续追踪的开放问题：

\begin{enumerate}
  \item 如何测量 long-horizon success，而不是只测单步 tool call？
  \item 如何在不无限保存隐私数据的前提下实现 personalization？
  \item 如何让 memory 同时支持时间、层级、关系和遗忘？
  \item 如何区分 reasoning trace 的解释价值与真实因果机制？
  \item 如何让 multi-agent verification 独立于原始错误通道？
  \item 如何定义可撤销、可补偿和不可逆 action 的不同执行策略？
  \item 如何让更小模型、专用模型与 retrieval 在能力和成本之间协作？
\end{enumerate}

\begin{warningbox}{最后的工程判断}
模型给出的 action proposal 必须与真实 side effect 分离。任何涉及资金、身份、发布、删除、权限或不可逆操作的系统，都应把模型放在受约束 control plane 后面；没有 confirmation、audit、recovery 和 sandbox 的 autonomy，不应被包装成可靠自动化。
\end{warningbox}

\subsection{自测问题}

\begin{enumerate}
  \item 为什么 attention 的 $\sqrt{d_k}$ scaling 有助于 softmax 稳定？
  \item self-attention 与 cross-attention 的 Q/K/V 来源有何不同？
  \item 为什么 emergent metric 可能在底层能力平滑时出现跳变？
  \item RLHF 与 DPO 分别在哪一步使用 preference data？
  \item MoE 的总参数量与 active parameters 为什么不同？
  \item RAG、fine-tuning、continual learning 与 model editing 各改变什么？
  \item CoT 的可读中间步骤为什么不必然是忠实解释？
  \item 单次 LLM call 缺少哪些 agent component？
  \item UI control 相比 API control 多了哪些风险？
  \item manager 为什么不能直接相信 worker 的 ``done''？
  \item 当 action 不可重复时，redo loop 应增加什么机制？
  \item LLM OS 类比中，哪些传统 OS 保证仍然缺失？
\end{enumerate}

\subsection{拓展阅读}

\begin{itemize}
  \item \href{https://arxiv.org/abs/1706.03762}{Vaswani et al., Attention Is All You Need}：Transformer 与 multi-head attention 的原始论文。
  \item \href{https://arxiv.org/abs/2206.07682}{Wei et al., Emergent Abilities of Large Language Models} 与 \href{https://arxiv.org/abs/2304.15004}{Schaeffer et al., Are Emergent Abilities of Large Language Models a Mirage?}：分别理解现象与 metric caveat。
  \item \href{https://arxiv.org/abs/2305.18290}{Rafailov et al., Direct Preference Optimization}：DPO objective 与 RLHF 关系。
  \item \href{https://arxiv.org/abs/2202.05262}{Meng et al., Locating and Editing Factual Associations in GPT}：ROME 与 factual model editing。
  \item \href{https://arxiv.org/abs/2201.11903}{Wei et al., Chain-of-Thought Prompting} 与 \href{https://arxiv.org/abs/2305.10601}{Yao et al., Tree of Thoughts}：线性 rationale 与搜索式 reasoning。
  \item \href{https://docs.google.com/presentation/d/1oXPs3LXtIVIsVbwTyGjAWj_aWvak9c1uNC4uhkS6glk/edit}{CS25 V4 official slide deck}：本讲 114 页 canonical source。
\end{itemize}

\end{document}
